\par\medskip
\hypertarget{native-079-article060}{}\pdfbookmark[2]{60.}{article-060}
\begin{center}60.\end{center}
Designantibus $x$, $y$, $z$ distantias corporis coelestis a tribus planis in Sole sub angulis rectis se secantibus; $X$, $Y$, $Z$ distantias terrae (sive centri sive puncti in superficie) ab iisdem planis; patet, $x-X$, $y-Y$, $z-Z$ fore distantias corporis coelestis a tribus planis illis parallele per terram ductis, hasque distantias ad distantiam corporis a terra ipsiusque \textit{locum geocentricum}\footnote{In sensu latiori: proprie enim haec expressio ad eum casum refertur, ubi recta e terrae \textit{centro} ducitur.}, i. e. situm proiectionis rectae a terra ad ipsum ductae in sphaera coelesti, relationem eandem habituras, quam $x$, $y$, $z$ habent ad distantiam a Sole locumque heliocentricum. Sit $\Delta$ distantia corporis coelestis a terra; concipiatur in sphaera \SourceLeaf{080}coelesti perpendiculum a loco geocentrico ad circulum maximum, qui respondet plano distantiarum $z$, demissum, sitque $a$ distantia intersectionis a polo positivo circuli maximi, qui respondet plano ipsarum $x$, denique sit $b$ longitudo ipsius perpendiculi sive distantia loci geocentrici a circulo maximo distantiis $z$ respondente. Tunc erit $b$ latitudo aut declinatio geocentrica, prout planum distantiarum $z$ est ecliptica aut aequator; contra $a+N$ longitudo seu ascensio recta geocentrica, si $N$ designat in casu priori longitudinem in posteriori ascensionem rectam poli plani distantiarum $x$. Quamobrem erit
\[
\begin{aligned}
x-X&=\Delta\cos b\cos a\\
y-Y&=\Delta\cos b\sin a\\
z-Z&=\Delta\sin b.
\end{aligned}
\]
Duae priores aequationes dabunt $a$ atque $\Delta\cos b$; quantitas posterior (quam positivam fieri oportet) cum aequatione tertia combinata dabit $b$ atque $\Delta$.
